\subsection{比变量}
\label{sec:app-b-specific-variables}

现在使用比变量, 即单位质量的量. 将系统的比能和比熵分别记为 $e$ 和 $s$. 此外, 引入比体积 $v=1/\rho$ (其中 $\rho$ 是所研究系统的密度) 作为体积变量. 上一节得到的关系变为
\[
 e=Ts-Pv,
\]
\[
 s\,dT-v\,dP=0.
\]

正如上面所看到的, 对于比体积恒定时的一次元可逆变换, 系统获得的热量写成 $\Delta Q=T\,ds$, 但同时还有 $\Delta Q=c_v\,dT$, 由此可推出
\[
 \left(\frac{\partial s}{\partial T}\right)_v=\frac{c_v}{T},
\]
这里用记号 $(\partial F/\partial x)_y$ 表示当变量 $y$ 取固定值时, 量 $F$ 关于 $x$ 的偏导数. 这个记号有助于明确每次计算中对独立状态变量的选择. 因而, 可以把比熵的微分 (把它视为 $T$ 和 $v$ 的函数) 写成
\[
 ds=\frac{c_v}{T}\,dT+\left(\frac{\partial s}{\partial v}\right)_T\,dv.
\]

由 Gibbs--Duhem 关系\eqref{eq:app-b-gibbs-duhem}, 有
\[
 d(Pv)=P\,dv+v\,dP=P\,dv+s\,dT;
\]
因此, 可以对任意 (光滑) 函数的交叉二阶导数使用 Schwarz 定理, 从而得到所谓 Maxwell 关系中的一个:
\[
 \left(\frac{\partial s}{\partial v}\right)_T
 =\left(\frac{\partial P}{\partial T}\right)_v.
\]
于是
\[
 ds=\frac{c_v}{T}\,dT+\left(\frac{\partial P}{\partial T}\right)_v\,dv,
\]
再回到比内能, 得到
\begin{equation}
 de=T\,ds-P\,dv
 =c_v\,dT+\left(T\left(\frac{\partial P}{\partial T}\right)_v-P\right)dv.
 \label{eq:app-b-specific-internal-energy}
\end{equation}

因此, 如果知道所研究流体的状态方程 (即 $P$ 关于 $T$ 和 $v$ 的表达式), 就可以计算最后这两个关系中的所有系数.

例如, 气体动力学中的一些经典状态方程如下:
\begin{itemize}
 \item 理想气体定律
 \[
  P=k\frac{T}{v},
 \]
 其中 $k$ 是一个只依赖于所研究流体摩尔质量的常数.
 \item Van der Waals 状态方程
 \[
  \left(P+\frac{a}{v^2}\right)\left(1-\frac{b}{v}\right)=k\frac{T}{v},
 \]
 其中 $a$ 和 $b$ 是另外两个依赖于所研究流体的常数. 与非常简单的理想气体状态方程不同, 这个定律考虑了分子的直径 (假定分子是硬球) 以及分子之间的吸引力.
\end{itemize}

液体也满足状态方程, 但在这种情形中, 正如 Chapter~\ref{chap:fluid-mechanics-equations} 的 Section~\ref{sec:chap1-incompressible-models} 所说明的, 流动通常可以视为不可压缩, 因而状态方程的具体表达式对于描述流动并无用处.

最后注意到, 在得到 Equation~\eqref{eq:app-b-specific-internal-energy} 时, 并未使用熵的显式表达式, 而且熵也没有出现在所得到的公式中. 仅仅假定这个状态变量 $S$ 存在, 就使我们能够应用 Schwarz 定理, 并由此推出基本的热力学关系. 在本书主要讨论的不可压缩框架中, 熵并不显式出现在描述流动的最终方程组中. 然而, 它是可压缩流动建模中的一个基本量.
